\documentclass[journal]{IEEEtran}

% PACKAGES
\usepackage{amsmath,amsfonts}
\usepackage{graphicx}
\usepackage{url}
\usepackage{booktabs} % For professional-looking tables
\usepackage{hyperref} % For clickable links (e.g., email)
\usepackage{balance}  % To balance the last page columns
\usepackage{makecell}  % for multi-line header cells
\usepackage{orcidlink}
\usepackage{flafter}
\usepackage{placeins}
\usepackage{float}
\usepackage{caption}

% PDF information and link colors
\hypersetup{
    pdftitle={QUANTITATIVE ANALYSIS OF THE IMPACT OF IMAGE PREPROCESSING ON THE ACCURACY OF COMPUTER VISION MODELS TRAINED FOR THE IDENTIFICATION OF CUTANEOUS DERMATOLOGICAL DISEASES},
    pdfauthor={Gabriel Mitelman Tkacz and Gustavo Scalabrini Sampaio},
    colorlinks=true,
    linkcolor=blue,
    citecolor=blue,
    urlcolor=blue,
}

% Correct bad hyphenation here
\hyphenation{op-tical net-works semi-conduc-tor}


\begin{document}

% ---
% TITLE AND AUTHOR INFORMATION
% ---
\title{Quantitative Analysis of the Impact of Image Preprocessing on the Accuracy of Computer Vision Models Trained for the Identification of Cutaneous Dermatological Diseases}

\author{Gabriel Mitelman Tkacz\orcidlink{0009-0004-3619-4561},~\IEEEmembership{Member,~IEEE,}
        and Gustavo Scalabrini Sampaio\orcidlink{0000-0003-1150-5584}% <-this % stops a space
\thanks{This work received funding from the Brazilian Federal Agency for Support and Evaluation of Graduate Education (CAPES Foundation).}% <-this % stops a space
\thanks{G. M. Tkacz and G. S. Sampaio are with the School of Computer Science and Informatics, Mackenzie Presbyterian University, São Paulo, SP, Brazil (e-mail: gabriel.tkacz@ieee.org; gustavo.sampaio@mackenzie.br).}}

% The paper headers
\markboth{IEEE Transactions on Image Processing,~Vol.~XX, No.~Y, Month~202X}%
{Tkacz \MakeLowercase{\textit{et al.}}: Quantitative Analysis of Image Preprocessing Impact}

% Make the title area
\maketitle

% ---
% ABSTRACT AND KEYWORDS
% ---
\begin{abstract}
This study aims to analyze the impact of image preprocessing on the accuracy of computer vision models for identifying skin diseases. The research is justified by the growing relevance of computer vision in the medical field, especially for the early diagnosis of skin diseases. The methodology consisted of collecting and preprocessing images of skin lesions, training and evaluating different computer vision models for image classification. It is expected that the study will contribute to the identification of preprocessing techniques that optimize the accuracy of computer vision models for the diagnosis of skin diseases, aiding in the creation of more accurate and efficient diagnostic tools. The results show that the careful selection of preprocessing steps maximizes the adjusted accuracy gain, while more complex pipelines or those that prioritize filtering excessively burden the process, compromising the final performance.
\end{abstract}

\begin{IEEEkeywords}
Artificial Intelligence, Computer Vision, Image Preprocessing, Machine Learning, Convolutional Neural Network, Medical Diagnosis, Dermatological Diagnosis, Computer-aided Detection (CADe), Computer-aided Diagnosis (CADx).
\end{IEEEkeywords}


% ---
% TEXTUAL ELEMENTS
% ---

% ----------------------------------------------------------
% Introduction
% ----------------------------------------------------------
\section{Introduction}
\IEEEPARstart{T}{he} field of dermatology has experienced significant adoption of AI-based tools, particularly for diagnosing skin diseases \cite{Krakowski2024}. Computer-aided diagnosis (CADx) can increase diagnostic accuracy for dermatological diseases, like melanomas, by up to 12\% \cite{10.1001/jamanetworkopen.2021.7249, Brinker2019}. This represents a substantial improvement over the 66.3\% baseline accuracy of human doctors \cite{Barnett2019}. However, AI model effectiveness in accurately identifying dermatological conditions depends heavily on input data quality, which typically consists of medical images \cite{ALHINAI20201, Brinker2019}.

Image preprocessing techniques, including normalization, resizing, and augmentation, play a crucial role in enhancing image quality and mitigating potential biases, thereby impacting AI model prediction accuracy and reliability \cite{RODRIGUES2020103542, 9548351, Tarawneh2024}. This study aims to conduct a quantitative analysis of the impact of various image preprocessing techniques on the accuracy of computer vision models trained to identify cutaneous dermatological diseases. By systematically evaluating the effects of different preprocessing methods, the study seeks to provide insights into the optimal strategies for preparing medical images, thereby improving the diagnostic performance of AI models in dermatology.

This research aims to contribute to developing more accurate AI models for dermatological diagnosis by identifying optimal image preprocessing pipelines, thereby establishing best practices that can streamline model development and deployment.

Enhanced accuracy and reliability of AI-based diagnostic tools can democratize healthcare by making specialized dermatological diagnosis accessible to broader populations, particularly in underserved areas. By optimizing diagnostic processes and reducing error potential, this research can ultimately improve patient outcomes and reduce healthcare costs.

The overall objective of this study, therefore, is to investigate the impact of image preprocessing on the accuracy of computer vision models aimed at identifying skin diseases. To achieve this, the specific goals include selecting and applying relevant image preprocessing techniques; constructing and training computer vision models using suitable public datasets; and conducting a comparative evaluation of the models' accuracy with and without image preprocessing.

% ----------------------------------------------------------
% Theoretical Framework
% ----------------------------------------------------------
\section{Theoretical Framework}

Specific dermatological applications of computer vision include precise skin lesion segmentation, morphological characteristic analysis, dimensional measurement and temporal evolution of lesions, and automatic classification of different skin diseases. Main challenges in dermatological contexts involve variability in image capture conditions across different clinics and hospitals, different camera types and lighting conditions, variations in patient skin color, presence of hair and other artifacts, and the need for efficient processing in clinical environments.

Fundamental preprocessing techniques include specialized edge detection algorithms for delineating skin lesions, advanced segmentation methods for isolating regions of interest, extraction of specific features for dermatological analysis (such as ABCDE patterns for melanomas), and deep learning techniques for classification. Computer vision provides essential tools and methods for automatic dermatological image analysis, serving as the means for extracting and analyzing relevant features \cite{ARORA2021102358}.

Relevant techniques include transfer learning using pre-trained models adapted for dermatology, data augmentation specific to dermatological images (including rotations, reflections, and color variations that preserve diagnostic features), and state-of-the-art architectures modified to capture specific skin lesion features. CNNs represent the state-of-the-art in medical image classification \cite{DASH2020106240}.

Image preprocessing, this article's main focus, comprises operations applied to raw dermatological images to improve their quality and standardize them for automatic analysis. In dermatology, these techniques are particularly important due to substantial variability in image capture conditions and patient characteristics. This is crucial in dermatological applications where image quality and standardization significantly impact accurate diagnosis, ranging from basic operations like resizing and normalization to specific techniques for hair removal and non-uniform lighting correction \cite{9780692}.

Main challenges in dermatological contexts include preserving diagnostic features during transformations (such as specific lesion colors and textures), standardizing processing protocols among different medical institutions, adapting to different skin types and colors, and handling common dermatological image artifacts (hair, bubbles, reflections).

For the purposes of this research, related scientific works were taken into consideration for the definition of the proposed methodology.

In the context of classifying HEp-2 cells for the diagnosis of autoimmune diseases \cite{RODRIGUES2020103542}, different image preprocessing strategies were evaluated before feeding them into CNNs. All images were resized to 224x224 pixels using bilinear interpolation. Six different preprocessing strategies were tested: (a) original image without processing, (b) mean subtraction, (c) contrast stretching, (d) contrast stretching with mean subtraction, (e) histogram equalization, and (f) histogram equalization with mean subtraction. Surprisingly, the best results were obtained using the original images without preprocessing, with the exception of ResNet-50 which performed better with contrast stretching. Inception-V3 achieved the highest accuracy (96.69\%) using original images without preprocessing, followed by AlexNet (94.33\%) and VGG-16 (92.63\%). When combined with data augmentation techniques, performance improved further, with AlexNet reaching 98.00\% accuracy on the original images.

In biometric systems for authentication based on dorsal hand veins (DHV) \cite{Tarawneh2024}, researchers evaluated different image preprocessing techniques' impact before CNN classification. The study used 500 images from 50 patients (10 images each, with 5 from right hand and 5 from left). To increase data volume, data augmentation techniques included rotation, resizing, cropping, and mirroring. Various preprocessing filters were tested, including mean and median filters for noise removal. Results showed that models without preprocessing achieved 70\% training set accuracy. Using mean filters for noise removal, accuracy increased significantly to 99\% under different training conditions. Mean filters proved superior to median filters, which achieved 96-98\% accuracy.

For automatic skin lesion segmentation in dermoscopic images \cite{ARORA2021102358}, researchers proposed an end-to-end deep learning framework based on (i) modified U-Net; (ii) incorporating Group Normalization (GN) instead of Batch Normalization; (iii) Attention Gates (AG) in skip connections to highlight relevant details; (iv) Tversky Loss (TL) as the objective function to balance precision and sensitivity. The study used four public databases ($PH^2$, ISIC 2016, ISIC 2017, and ISIC 2018), totaling thousands of annotated images for training and testing. To increase model robustness, image preprocessing and data augmentation techniques were applied, including rotation, resizing, cropping, and mirroring, aiming to extract lesion contours more precisely. While original U-Net achieved limited results, the proposed model reached accuracies of 96\% on $PH^2$, 97\% on ISIC 2016, 95\% on ISIC 2017, and 96\% on ISIC 2018, outperforming reference methods on all evaluated databases.

For CADx systems assessing psoriasis severity \cite{DASH2020106240}, researchers developed a fully automated deep learning framework based on CNNs, organized into three integrated modules: (i) recognizing images as psoriasis or non-psoriasis; (ii) automatic segmentation of psoriatic lesions using modified U-Net; and (iii) assessing lesion severity using VGG-16 adapted to extract specific discriminative features. The study used extensive dermoscopic images from psoriasis patients, evaluated through k-fold cross-validation to ensure model robustness and generalization capability. To enhance training, image preprocessing (intensity normalization and artifact removal) and data augmentation techniques, including rotation, cropping, resizing, and mirroring, were applied. Compared to traditional machine learning methods, the cascaded system showed superior performance in all stages: high accuracy in lesion detection, high Dice and IoU scores in segmentation, and consistent accuracy in classifying severity levels, maintaining stable performance across different subset sizes.

% ----------------------------------------------------------
% Methodology
% ----------------------------------------------------------
\section{Methodology}

The proposed process constituted a quantitative analysis of the impact of different image preprocessing techniques on the accuracy of computer vision models trained to identify dermatological diseases. The methodology employed a binary CNN architecture for classifying skin images as healthy or diseased \cite{DASH2020106240}, implemented using the PyTorch framework with a subset of raw data from the HAM10000 dataset \cite{DBLP:journals/corr/abs-1803-10417}.

The implemented CNN architecture consisted of four sequential convolutional blocks, each composed of a convolutional layer, batch normalization, ReLU activation function, and max pooling. The first block started with 32 filters, progressively doubling to 256 filters in the fourth block. The network concluded with fully connected layers that reduced dimensionality to a single output with sigmoid activation, appropriate for binary classification. The dataset was managed through a custom class, organizing images into two categories: healthy skin and diseased skin.

The experimentation process was structured in hierarchical classes, starting with a base model (Class 0) without preprocessing, followed by models with single transformation (Class 1), two transformations (Class 2), three transformations (Class 3), and four transformations (Class 4). Investigated transformations included: normalization (N), equalization (E), denoising (D), and color space transformation (CS) as specified in Table~\ref{tab:params} \cite{Brinker2019, RODRIGUES2020103542, Tarawneh2024, computers8030062}.

% \begin{figure}[!ht]
%  \centering
%  \includegraphics[width=\columnwidth]{Group_1.png}
%  \caption{Quantitative analysis of preprocessing.}
%  \label{fig:quant}
% \end{figure}

The selection of optimal preprocessing parameters was conducted through systematic grid search optimization across each transformation type. For color space transformation, four target spaces were evaluated (BGR, HSV, LAB, and YUV). The denoising operation parameters were optimized by testing template window sizes ranging from 5 to 10 pixels and search window sizes from 20 to 25 pixels, yielding 36 unique combinations. For normalization parameters, a comprehensive grid search evaluated mean values and standard deviations ranging from 0.15 to 0.90 in increments of 0.15, generating 36 distinct parameter combinations. This parameter optimization process was conducted on the same dataset split used for the main experiments, ensuring consistency in evaluation conditions. Each parameter combination underwent complete model training and evaluation, with precision metrics and processing times recorded to inform the final parameter selection presented in Table~\ref{tab:analysis}.

Normalization reduces lighting and contrast differences between samples, allowing neural network optimizers to converge more quickly and stably, minimizing scale bias between input attributes. In skin lesion classification tasks, normalization helps networks focus on relevant textural and chromatic patterns rather than absolute intensity fluctuations \cite{irjet}.

Histogram equalization highlights subtle contour and texture details in skin lesions, enhancing discriminating feature extraction by convolutional filters \cite{10.5555/1076432, 10810559, 9025640}.

High-frequency noise removal suppresses thermal or acquisition artifacts without compromising relevant edges and structures. Non-Local Means-based algorithms assess similarity between small image regions to perform weighted averaging, attenuating isolated noise points while preserving repetitive patterns typical of dermatological lesions \cite{MAFI2019236}. This filtering improves signal-to-noise ratio, favoring model generalization on noisy data or data acquired under varied clinical conditions \cite{elad2023imagedenoisingdeeplearning}.

Converting images between different color spaces (e.g., RGB $\to$ HSV, RGB $\to$ LAB, or RGB $\to$ YUV) allows separation of luminance and chrominance information or exploration of more perceptually uniform representations \cite{WANG2018157}. In HSV space, for example, the hue channel can highlight color variations in skin rashes, while LAB space approximates human color perception \cite{ballester2022influencecolorspacesdeep}. These transformations reveal attributes complementary to those offered by RGB space, broadening the range of patterns captured by CNNs in detecting psoriasis and other skin diseases.

To comprehensively evaluate preprocessing impact across different performance regimes, experiments were conducted under varying baseline conditions. This approach revealed that baseline model accuracy serves as a critical modulating factor for preprocessing effectiveness. The relationship between baseline performance and preprocessing benefit follows a non-linear pattern, where higher baseline accuracies exponentially reduce the magnitude of observable improvements while simultaneously dampening the severity of performance degradation from suboptimal preprocessing choices.

The evaluation methodology divided the dataset into three subsets: training (80\%), testing (10\%), and validation (10\%), with a fixed seed to maintain consistent randomness across various models. The training process used the Adam optimizer with a learning rate of 0.0001 and the Mean Squared Error (MSE) loss function. Transformation combinations were systematically evaluated using permutations of available transformations, allowing comprehensive analysis of different preprocessing sequence impacts.

The results were presented in terms of three factors.

The $\alpha$ factor is defined as the absolute difference between the evaluated preprocessing model's accuracy ($\epsilon_x$) and the original model's accuracy ($\epsilon_0$):
\begin{equation}
\alpha = \epsilon_x - \epsilon_0
\end{equation}

Similarly, the $\gamma$ factor is defined as the ratio between the evaluated preprocessing model's training time ($\rho_x$) and the original model's training time ($\rho_0$):
\begin{equation}
\gamma = \frac{\rho_x}{\rho_0}
\end{equation}

To integrate computational cost into the accuracy gain metric, the weighted alpha ($\alpha_w$) factor is defined by:
\begin{equation}
\alpha_w = f(\alpha, \gamma) =
\begin{cases}
\frac{\alpha}{\gamma}, & \text{if } \alpha \geq 0 \\
\alpha \cdot \gamma, & \text{if } \alpha < 0
\end{cases}
\label{eq:1}
\end{equation}

Thus, combinations that increase accuracy with low time cost receive higher $\alpha_w$ values, while those that penalize training time have their gain adjusted according to factor $\gamma$. An important methodological consideration concerns the relationship between baseline model accuracy and the magnitude of potential improvements from preprocessing. When the baseline accuracy ($\epsilon_0$) approaches the theoretical maximum of 100\%, the available margin for improvement becomes increasingly constrained, following a diminishing returns pattern. This phenomenon, known as the ceiling effect \cite{vogt2005scale}, fundamentally alters the interpretation of preprocessing benefits. For instance, with a $\epsilon_{0}$ of 80\%, the maximum possible improvement ($\alpha_{max}$) is 20 percentage points, whereas a $\epsilon_{0}$ of 99\% restricts $\alpha_{max}$ to merely 1. This compression of the improvement space necessitates careful consideration when comparing preprocessing effectiveness across different baseline conditions. The relationship can be expressed as the maximum potential gain $\alpha_{max} = 100 - \epsilon_0$, which decreases linearly as baseline performance improves. Implementation included reproducibility mechanisms through random seed control and used CUDA GPU acceleration when available. The experimentation process was automated and instrumented to collect detailed metrics during training and evaluation, including training loss, validation loss, accuracy, confusion matrices, and execution time. Experiments were performed on a system equipped with an Intel Core i5-12600KF processor (10 cores, 16 threads, 3.7 GHz), 32 GB of 3200 MHz RAM, and an NVIDIA GeForce RTX 4080 SUPER GPU running CUDA 12.0.140.


% ----------------------------------------------------------
% Results
% ----------------------------------------------------------
\section{Results}

% Using table* to span two columns for the large table
\begin{table*}[h!t]
\centering
\caption{Model metrics after preprocessing}
\label{tab:analysis}
\begin{tabular*}{\textwidth}{@{\extracolsep{\fill}} l l l l r r r r @{}}
\toprule
\textbf{transform\_1} &
\textbf{transform\_2} &
\textbf{transform\_3} &
\textbf{transform\_4} &
\makecell{\boldsymbol{
\alpha} \\ $(\epsilon_0 = 88.5)$} &
\makecell{\boldsymbol{\alpha_w} \\ $(\epsilon_0 = 88.5)$} &
\makecell{\boldsymbol{\alpha} \\ $(\epsilon_0 = 98.3)$} &
\makecell{\boldsymbol{\alpha_w} \\ $(\epsilon_0 = 98.3)$} \\
\midrule
CS & - & - & - & 0.50 & 0.69 & 0.00 & 0.00 \\
CS & D & - & - & 2.50 & 0.85 & 0.40 & 0.06 \\
CS & D & E & - & 2.50 & 0.84 & 0.00 & 0.00 \\
CS & D & E & N & 7.50 & 2.50 & 0.00 & 0.00 \\
CS & D & N & - & 7.50 & 2.56 & -0.20 & -0.03 \\
CS & D & N & E & -28.00 & -85.10 & -0.20 & -0.03 \\
CS & E & - & - & 2.00 & 2.70 & -0.30 & -0.27 \\
CS & E & D & - & 0.00 & 0.00 & 0.00 & 0.00 \\
CS & E & D & N & 6.50 & 2.12 & 0.00 & 0.00 \\
CS & E & N & - & 9.00 & 12.17 & -0.40 & -0.41 \\
CS & E & N & D & -15.00 & -45.96 & -0.40 & -0.41 \\
CS & N & - & - & 7.00 & 7.58 & 0.30 & 0.29 \\
CS & N & D & - & -4.00 & -11.95 & -0.50 & -0.07 \\
CS & N & D & E & -1.50 & -4.56 & -0.50 & -0.07 \\
CS & N & E & - & -6.50 & -4.79 & -0.20 & -0.18 \\
CS & N & E & D & -2.50 & -7.66 & -0.20 & -0.18 \\
D & - & - & - & 8.00 & 2.75 & 0.40 & 0.06 \\
D & CS & - & - & -0.50 & -1.38 & 0.40 & 0.06 \\
D & CS & E & - & 0.50 & 0.17 & 0.30 & 0.04 \\
D & CS & E & N & 7.50 & 2.51 & 0.30 & 0.04 \\
D & CS & N & - & 5.00 & 1.71 & 0.10 & 0.01 \\
D & CS & N & E & -5.00 & -14.95 & 0.10 & 0.01 \\
D & E & - & - & 6.00 & 2.20 & -0.80 & -0.12 \\
D & E & CS & - & 0.50 & 0.17 & -0.10 & -0.01 \\
D & E & CS & N & 8.00 & 2.96 & -0.10 & -0.01 \\
D & E & N & - & 9.50 & 3.21 & -0.70 & -0.10 \\
D & E & N & CS & -15.50 & -42.22 & -0.70 & -0.10 \\
D & N & - & - & 8.50 & 2.89 & 0.40 & 0.06 \\
D & N & CS & - & -27.00 & -79.73 & -0.80 & -0.12 \\
D & N & CS & E & -29.50 & -79.30 & -0.80 & -0.12 \\
D & N & E & - & -2.00 & -5.93 & -0.90 & -0.13 \\
D & N & E & CS & -27.00 & -80.62 & -0.90 & -0.13 \\
E & - & - & - & 8.00 & 10.97 & -0.30 & -0.29 \\
E & CS & - & - & 6.50 & 6.85 & -0.60 & -0.60 \\
E & CS & D & - & 3.50 & 1.17 & -0.10 & -0.01 \\
E & CS & D & N & 9.50 & 3.14 & -0.10 & -0.01 \\
E & CS & N & - & 9.50 & 13.17 & -0.40 & -0.37 \\
E & CS & N & D & -11.00 & -33.77 & -0.40 & -0.37 \\
E & D & - & - & 9.00 & 3.32 & -1.30 & -0.19 \\
E & D & CS & - & 6.00 & 1.99 & -1.00 & -0.15 \\
E & D & CS & N & 8.50 & 2.78 & -1.00 & -0.15 \\
E & D & N & - & 8.50 & 2.86 & -1.70 & -0.25 \\
E & D & N & CS & -8.50 & -24.71 & -1.70 & -0.25 \\
E & N & - & - & 8.00 & 11.01 & -0.50 & -0.47 \\
E & N & CS & - & 3.00 & 4.00 & -1.50 & -1.35 \\
E & N & CS & D & 2.50 & 0.82 & -1.50 & -1.35 \\
E & N & D & - & 6.50 & 2.19 & -3.50 & -0.51 \\
E & N & D & CS & -3.00 & -9.19 & -3.50 & -0.51 \\
N & - & - & - & 7.00 & 9.99 & 0.00 & 0.00 \\
N & CS & - & - & -8.50 & -5.95 & -0.50 & -0.53 \\
N & CS & D & - & -8.00 & -24.08 & -0.80 & -0.12 \\
N & CS & D & E & -6.50 & -19.47 & -0.80 & -0.12 \\
N & CS & E & - & -10.00 & -7.12 & -0.80 & -0.72 \\
N & CS & E & D & -7.00 & -21.35 & -0.80 & -0.72 \\
N & D & - & - & -4.00 & -11.75 & -0.80 & -0.12 \\
N & D & CS & - & -4.00 & -12.34 & -0.50 & -0.07 \\
N & D & CS & E & -19.50 & -59.47 & -0.50 & -0.07 \\
N & D & E & - & -2.50 & -7.45 & -0.70 & -0.10 \\
N & D & E & CS & -23.50 & -64.23 & -0.70 & -0.10 \\
N & E & - & - & -2.50 & -1.85 & -1.10 & -1.04 \\
N & E & CS & - & -8.50 & -6.05 & -0.70 & -0.64 \\
N & E & CS & D & -16.00 & -43.78 & -0.70 & -0.64 \\
N & E & D & - & -2.00 & -6.03 & -0.60 & -0.09 \\
N & E & D & CS & -11.50 & -35.06 & -0.60 & -0.09 \\
\bottomrule
\end{tabular*}
\end{table*}

Table~\ref{tab:analysis} presents the $\alpha$ and $\alpha_w$ values obtained for different transformation sequences applied in image preprocessing. Columns \emph{transform\_1} to \emph{transform\_4} indicate the execution order of operations. The value of $\alpha$ reflects gross accuracy gain, while $\alpha_w$ adjusts this gain according to associated computational cost, emphasizing or attenuating the absolute value of $\alpha$ based on model training time, as shown in \eqref{eq:1}.

A critical observation emerges when examining preprocessing impact under different baseline conditions. When the baseline model achieves substantially higher initial accuracy, the entire distribution of $\alpha$ values undergoes dramatic compression. This compression effect is non-uniform across the performance spectrum, with positive gains experiencing more severe attenuation than negative impacts. The phenomenon suggests that preprocessing techniques may have fundamentally different effects depending on the intrinsic difficulty of the classification task and the baseline model's capacity to extract relevant features from unprocessed images.

\begin{figure}[H]
  \centering
  \includegraphics[width=\columnwidth]{top5.png}
  \caption{Visual comparison of a sample input (left) versus the results of the top 5 performing preprocessing pipelines. This qualitative analysis highlights how specific feature extraction is facilitated by the transformation sequences listed in Table~\ref{tab:analysis}.}
  \label{fig:top5_pipelines}
\end{figure}

The experimental results reveal complex relationships between preprocessing sequences, accuracy gains, and computational costs that challenge conventional assumptions about image preprocessing benefits in medical imaging applications.

Analysis demonstrates that optimal preprocessing strategies achieve $\alpha$ of up to 9.50\% on a model with suboptimal $\epsilon_0$. Three transformation sequences reached this maximum gain, though with markedly different computational efficiency profiles. The sequence combining equalization, color space transformation, and normalization (E → CS → N) emerged as the most effective overall strategy, achieving the maximum accuracy gain while maintaining an $\alpha_w$ of 13.17, indicating exceptional computational efficiency. In contrast, the sequences incorporating denoising operations (D → E → N with $\alpha_w$ = 3.21, and E → CS → D → N with $\alpha_w$ = 3.14) achieved identical accuracy improvements but at substantially higher computational costs, reducing their practical utility when cost-adjusted performance is considered.

The relationship between preprocessing complexity and performance exhibits notable non-linearity. Single-operation preprocessing strategies demonstrate surprisingly competitive performance when computational costs are considered. For example, isolated equalization achieved an $\alpha$ of 8.00\% and $\alpha_w$ of 10.97, while normalization alone yielded a $\alpha$ of 7.00\% with an $\alpha_w$ of 9.99. These findings suggest that simpler preprocessing approaches may offer superior cost-benefit ratios compared to complex multi-stage pipelines. This pattern becomes particularly evident when examining sequences with three or four transformations, where moderate $\alpha$ values of 6.50\% to 7.50\% are offset by computational burdens that reduce $\alpha_w$ values to the range of 1.99 to 2.78, representing a four to five-fold reduction in cost-adjusted performance.

\begin{table}[H]
\centering
\caption{Parameters used for preprocessing}
\label{tab:params}
\begin{tabular}{lll}
\toprule
Preprocessing & Parameters (low $\epsilon_{0}$ — high $\epsilon_{0}$) \\
\midrule
N & \begin{tabular}[t]{@{}l@{}}Mean: 0.4 — 0.15\\ Standard deviation: 0.2 — 0.75\end{tabular} \\
E & - \\
D & \begin{tabular}[t]{@{}l@{}}Model window size: 5 — 9\\ Search window size: 19 — 22\end{tabular} \\
CS & RGB2HSV — RGB2LAB \\
\bottomrule
\end{tabular}
\end{table}

A critical finding emerges from the substantial proportion of preprocessing sequences that actively degraded model performance. Approximately half of all tested combinations produced negative $\alpha$ values, with the most detrimental sequence (D → N → CS → E) reducing accuracy by 29.50 percentage points and yielding an $\alpha_w$ of -79.30. The analysis reveals systematic patterns in these failures, with sequences beginning with normalization followed by color space conversion consistently underperforming, suggesting fundamental incompatibilities between certain transformation orderings and the network's feature extraction mechanisms. The computational penalty applied to negative $\alpha$ values through \eqref{eq:1} further amplifies these poor results, with the worst-performing sequences reaching $\alpha_w$ values as low as -85.10.

\begin{figure}[H]
  \centering
  % Subfigure 1: Low Baseline
  \begin{minipage}[b]{0.8\columnwidth}
    \centering
    \includegraphics[width=\columnwidth]{norm.png}
    \captionsetup{labelformat=empty}
    \caption {a) Low Baseline ($\epsilon_0 = 88.5\%$)}
  \end{minipage}
  \vfill
  % Subfigure 2: High Baseline
  \begin{minipage}[b]{0.8\columnwidth}
    \centering
    \includegraphics[width=\columnwidth]{norm_high_con.png}
    \captionsetup{labelformat=empty}
    \caption {b) High Baseline ($\epsilon_0 = 98.3\%$)}
  \end{minipage}
  \caption{Comparison of Normalization (N) parameter correlations under different baseline conditions.}
  \label{fig:norm-comparison}
\end{figure}

The sensitivity of these findings to baseline model performance represents a crucial consideration for practical implementation. Supplementary experiments conducted with a substantially higher $\epsilon_0$ (and thus a much lower $\alpha_{max}$) revealed dramatic compression of the entire performance distribution. Under these high-baseline conditions, the maximum achievable $\alpha$ decreased to merely 0.40, a 24-fold reduction compared to the original experimental setting. This compression effect manifests asymmetrically, with positive gains experiencing more severe attenuation than negative impacts. The most detrimental preprocessing combination under high-baseline conditions produced only a $\alpha$ of -4.80, compared to -29.50 observed at lower baseline accuracy, representing a six-fold reduction in negative impact.

This baseline-dependent behavior fundamentally alters the interpretation of preprocessing effectiveness. When models operate near their theoretical performance ceiling, the marginal utility of preprocessing diminishes exponentially while computational costs remain constant or even increase due to the additional processing overhead. The phenomenon suggests that preprocessing benefits cannot be considered universal constants but rather context-dependent variables whose magnitude correlates inversely with baseline model performance. In practical terms, this implies that preprocessing strategies optimized for challenging datasets or lower-performing models may provide negligible benefits when applied to high-quality imaging scenarios or well-tuned architectures.

The temporal cost analysis reveals that denoising operations consistently impose the highest computational burden, with time ratios exceeding 6.5 times the baseline training duration. This substantial overhead explains why sequences incorporating denoising early in the pipeline consistently underperform in terms of $\alpha_w$, despite sometimes achieving respectable raw accuracy gains. Conversely, normalization and color space transformations typically maintain time ratios near unity, explaining their superior cost-adjusted performance metrics.

\begin{figure}[H]
  \centering
  % Subfigure 1: Low Baseline
  \begin{minipage}[b]{0.7\columnwidth}
    \centering
    \includegraphics[width=\columnwidth]{cs.png}
    \captionsetup{labelformat=empty}
    \caption {a) Low Baseline ($\epsilon_0 = 88.5\%$)}
  \end{minipage}
  \vfill
  % Subfigure 2: High Baseline
  \begin{minipage}[b]{0.7\columnwidth}
    \centering
    \includegraphics[width=\columnwidth]{cs_high_con.png}
    \captionsetup{labelformat=empty}
    \caption {b) High Baseline ($\epsilon_0 = 98.3\%$)}
  \end{minipage}
  \caption{Comparison of Color Space (CS) parameter correlations under different baseline conditions.}
  \label{fig:cs-comparison}
\end{figure}

These findings collectively indicate that effective preprocessing strategy selection requires careful consideration of multiple factors including baseline model performance, computational resources, and the specific characteristics of the target dataset. The results strongly support adopting simpler preprocessing pipelines, particularly in scenarios with reasonable baseline performance, while reserving complex multi-stage preprocessing for situations where baseline accuracy is substantially below acceptable thresholds and computational resources are not constraining factors.

% ----------------------------------------------------------
% Conclusion
% ----------------------------------------------------------
\section{Discussion}

Comprehensive investigation into the impact of image preprocessing on computer vision models for dermatological disease identification yields fundamental insights that challenge prevailing assumptions about preprocessing benefits in medical imaging applications. Through systematic evaluation of 64 distinct preprocessing pipelines applied to a binary classification task using the HAM10000 dataset, it was demonstrated that the relationship between preprocessing techniques and model performance exhibits far greater complexity than previously recognized. Findings reveal that while carefully selected preprocessing sequences can yield substantial accuracy improvements of up to 9.5 percentage points in primary experiments, an equal proportion of preprocessing combinations actively degrade model performance, with the most detrimental sequences reducing accuracy by nearly 30 percentage points.

The introduction of a weighted accuracy metric ($\alpha_w$) helps to address a critical variable in preprocessing evaluation frameworks. Traditional approaches have predominantly focused on raw accuracy gains without adequately considering the computational overhead imposed by preprocessing operations. Analysis demonstrates that this oversight can lead to misleading conclusions about preprocessing effectiveness. For instance, while three distinct preprocessing sequences achieved the maximum accuracy gain of 9.5\%, their $\alpha_w$ values ranged from 3.14 to 13.17, representing a four-fold difference in practical utility when computational costs are considered. This disparity becomes particularly pronounced for denoising operations, which consistently imposed computational burdens exceeding six times the baseline training duration, effectively negating their accuracy benefits in resource-constrained clinical environments.

Perhaps the most significant revelation emerging from this research concerns the profound influence of baseline model performance on preprocessing effectiveness. A supplementary analysis, conducted with a baseline accuracy of 98.3\% compared to the initial 88.5\%, exposed a dramatic compression of the entire performance distribution that fundamentally alters the cost-benefit calculus of preprocessing implementation. The maximum achievable improvement decreased by a factor of 24, from 9.5\% to merely 0.4\%, while the computational costs remained constant or increased. This finding has immediate practical implications for clinical deployment strategies, suggesting that resources allocated to preprocessing optimization in high-performance scenarios might be better directed toward data acquisition improvements or model architecture refinement.

The observed patterns of preprocessing failure provide valuable insights into the underlying mechanisms through which these transformations interact with neural network feature extraction processes. Results indicate that certain transformation orderings, particularly those beginning with normalization followed by color space conversion, consistently interfere with the network's ability to extract diagnostically relevant features. This systematic failure pattern suggests that preprocessing transformations are not merely independent image enhancement operations but rather coupled processes that fundamentally alter the feature space in ways that can either facilitate or impede subsequent learning. The asymmetric nature of these effects, where beneficial combinations cluster around specific transformation sequences while detrimental combinations are more broadly distributed, implies that successful preprocessing requires understanding not just individual transformation effects but their synergistic interactions within the complete pipeline.

The implications of these findings extend beyond the immediate context of dermatological imaging to broader questions about transfer learning and model generalization in medical AI applications. The extreme sensitivity of preprocessing benefits to baseline performance suggests that preprocessing strategies developed and validated on one dataset or model architecture may not transfer effectively to different contexts. This challenges the practice of adopting preprocessing pipelines from published literature without careful validation in the specific deployment context. Furthermore, the inverse relationship between baseline performance and preprocessing benefit indicates that as medical imaging AI systems continue to improve, the utility of complex preprocessing pipelines will diminish, potentially reaching a point where preprocessing becomes counterproductive except in specialized scenarios involving particularly challenging image conditions or a lack of data.

From a clinical deployment perspective, findings support a paradigm shift toward adaptive preprocessing strategies that dynamically adjust based on image quality assessment and model confidence metrics. Rather than applying uniform preprocessing pipelines to all images, future systems should implement conditional preprocessing that reserves computationally intensive transformations for images where baseline model confidence falls below acceptable thresholds. This approach would optimize the trade-off between computational efficiency and diagnostic accuracy while maintaining robust performance across the heterogeneous image quality conditions typical of real-world clinical practice. The development of such adaptive systems represents a natural evolution from the static preprocessing pipelines currently prevalent in medical imaging applications.

The methodological limitations of this study, while providing opportunities for future research, also highlight important considerations for the broader field. The focus on four fundamental preprocessing techniques, while enabling systematic exploration of their combinations, necessarily excludes more sophisticated preprocessing approaches such as learned preprocessing layers, attention-guided enhancement, or domain-specific artifact removal techniques. The restriction to a single CNN architecture and binary classification task, though enabling controlled comparison, limits the generalizability of specific preprocessing recommendations. Future investigations should explore how preprocessing effectiveness varies across different neural network architectures, particularly modern architectures incorporating attention mechanisms or vision transformers, which may exhibit fundamentally different sensitivities to preprocessing transformations.

Looking toward future research directions, several critical questions emerge from these findings. The mechanism through which certain preprocessing sequences interfere with feature extraction warrants detailed investigation through visualization and interpretation techniques that can illuminate the interaction between preprocessing transformations and learned feature representations such as Grad-CAM \cite{8237336}. The development of theoretical frameworks for predicting preprocessing effectiveness based on dataset characteristics and model architecture would provide valuable guidance for practitioners. Additionally, the exploration of learnable preprocessing modules that can adapt to specific datasets and tasks represents a promising direction for achieving optimal preprocessing without extensive manual experimentation.


% ----------------------------------------------------------
% Conclusion
% ----------------------------------------------------------
\section{Conclusion}

This research demonstrates that the role of preprocessing in medical image analysis is far more nuanced than conventionally assumed. While preprocessing can provide substantial benefits under appropriate conditions, its effectiveness is strongly modulated by baseline model performance, computational constraints, and the specific ordering of transformations. The evidence presented here advocates for a more thoughtful, context-aware approach to preprocessing implementation that considers not just potential accuracy gains but also computational costs, baseline performance levels, and the specific characteristics of the deployment environment. As the field of medical AI continues to mature, the insights provided by this research will hopefully become increasingly relevant for developing efficient, scalable, and clinically viable diagnostic systems. The ultimate goal of improving patient outcomes through AI-assisted diagnosis requires not just accurate models but also practical, deployable solutions that balance performance with real-world constraints.

% ----------------------------------------------------------
% Bibliographical references
% NOTE: The original .bib file was not provided.
% ----------------------------------------------------------
% Set the bibliography style to the IEEE standard
\bibliographystyle{IEEEtran}

\bibliography{refs}

\balance

\end{document}